% !PRINTED 335
\weberpage{335}
Die linearen Faktoren von $N(\lambda)$ und $N(\lambda')$ müssen also, von einem konstanten Faktor abgesehen, miteinander übereinstimmen, weil man eine ganze Funktion beliebiger Variablen nur auf eine Weise in irreducible (hier lineare) Faktoren zerlegen kann (Bd.\,I, § 20, Bd.\,II, § 152), und da alle diese Faktoren demselben Körper $\Omega$ angehören, so gibt es unter den konjugierten Faktoren $\lambda'$ einen, der der Bedingung genügt:
\[
\lambda = \eta\lambda',
\]
worin $\eta$ eine Zahl in $\Omega$ ist. Die Funktionale $\lambda$ und $\lambda'$ und demnach die entsprechenden Ideale sind also äquivalent.

4. \gesperrt{Gehört in einem Normalkörper zu den zwei Idealen $\mathfrak{a}$ und $\mathfrak{a}'$ dieselbe Formenklasse, so gibt es unter den mit $\mathfrak{a}'$ konjugierten Idealen eines, das mit $\mathfrak{a}$ äquivalent ist.}

\begin{center}
\textbf{§ 89. Komposition der Formen und Multiplikation der Ideale.}
\end{center}

Es seien jetzt $\mathfrak{a}$ und $\mathfrak{b}$ zwei Ideale in $\Omega$ und
\weq{1}{\mathfrak{c} = \mathfrak{a}\mathfrak{b}}
das aus beiden gebildete Produkt. Es seien ferner
\weq{2}{\begin{aligned}
\mathfrak{a} &= (\alpha_1,\ \alpha_2,\ \ldots,\ \alpha_n), \\
\mathfrak{b} &= (\beta_1,\ \beta_2,\ \ldots,\ \beta_n), \\
\mathfrak{c} &= (\gamma_1,\ \gamma_2,\ \ldots,\ \gamma_n)
\end{aligned}}
Basen dieser Ideale. Da jede Zahl $\alpha_h\beta_k$ zu $\mathfrak{c}$ gehört, so gibt es ein System ganzer rationaler Zahlen $c^{(r)}_{hk}$, so daß
\weq{3}{\alpha_h\beta_k = \overset{r}{\Sigma}\, c^{(r)}_{h,k}\,\gamma_r}
wird, und da jede Zahl in $\mathfrak{c}$ durch Multiplikation je einer Zahl aus $\mathfrak{a}$ und aus $\mathfrak{b}$ und Addition dieser Produkte entsteht (Bd.\,II, § 169), so ist auch umgekehrt
\weq{4}{\gamma_s = \overset{h,k}{\Sigma}\, a^{(s)}_{h,k}\,\alpha_h\beta_k,}
worin die $a^{(s)}_{h,k}$ gleichfalls ganze rationale Zahlen sind.

Setzt man (3) in (4) ein, so folgt
\weq{5}{\gamma_s = \overset{r}{\Sigma}\,\gamma_r \overset{h,k}{\Sigma}\, a^{(s)}_{h,k}\,c^{(r)}_{h,k},}
und daraus
\weq{6}{\overset{h,k}{\Sigma}\, a^{(s)}_{h,k}\,c^{(r)}_{h,k} = (r,s),}
worin
\[
(r,r) = 1, \quad (r,s) = 0, \quad r \neq s.
\]
